\chapter{Hiện thực và kiểm thử}
\label{chap:hienthuc}

% TODO: toàn bộ chương này cần nhóm điền nội dung thực tế từ quá trình phát triển.

\section{Hiện thực}
\label{sec:hienthuc}

% TODO: toàn bộ chương này cần nhóm điền nội dung thực tế từ quá trình phát triển. Khung mục dưới đây bám theo Ch.3 để đảm bảo tính nhất quán thiết kế--hiện thực.

\subsection{Môi trường và công cụ phát triển}
\label{sec:moitruong}
Mã nguồn được quản lý trên GitHub. Mỗi lần đẩy mã, GitHub Actions tự chạy kiểm tra kiểu, phân tích tĩnh bằng ESLint và kiểm thử kèm đo độ phủ bằng Jest; cổng độ phủ chặn mã không đạt ngưỡng tối thiểu của phần nghiệp vụ lõi trước khi hợp nhất. Kiểm thử tải bằng Grafana k6 và đo giao diện bằng Lighthouse chạy riêng theo đợt vì tốn thời gian.

Công nghệ dùng khi hiện thực bám đúng các lựa chọn tại mục~\ref{sec:giaiphapcongnghe}:
\begin{itemize}
    \item \textbf{Kho mã:} một monorepo quản lý bằng pnpm workspaces và Turborepo, gồm ứng dụng mua hàng, ứng dụng nội bộ, máy chủ nghiệp vụ và các gói dùng chung (kiểu dữ liệu TypeScript, thư viện thành phần giao diện).
    \item \textbf{Giao diện:} ứng dụng mua hàng bằng Next.js, đóng gói dạng PWA; ứng dụng nội bộ bằng React, Vite và Refine.
    \item \textbf{Máy chủ nghiệp vụ:} NestJS trên Node.js và TypeScript; truy cập dữ liệu qua Prisma, riêng khóa dòng tồn viết bằng SQL thô.
    \item \textbf{Dữ liệu:} PostgreSQL kèm các tiện ích pgvector, \texttt{unaccent}, \texttt{pg\_trgm}; Redis cho hàng đợi sự kiện và cache; MinIO cho tệp.
    \item \textbf{Dịch vụ AI:} Python với FastAPI; mô hình nhúng BGE-M3; thư viện dự báo chuỗi thời gian cho SARIMA, Prophet.
    \item \textbf{Bên thứ ba:} cổng thanh toán và hóa đơn điện tử ở dạng bản giả lập đặt sau giao diện tích hợp; email qua SMTP. Kết nối VNPay, Giao Hàng Nhanh ở môi trường thử nghiệm và dịch vụ đẩy của trình duyệt (Web Push) \emph{(chỉ thiết kế)}.
\end{itemize}

% TODO: ghi phiên bản cụ thể của từng thành phần khi dựng xong môi trường; ngưỡng cổng chất lượng đã đặt.

\subsection{Kiến trúc triển khai}
\label{sec:trienkhai}
% TODO: mô tả hạ tầng triển khai thực tế (server/cloud, container hóa nếu có), đối chiếu với kiến trúc thiết kế tại mục~\ref{sec:kientruc}.

\subsection{Hiện thực các module chính}
\label{sec:hienthucmodule}

Dựa trên kiến trúc và thiết kế chi tiết đã đề xuất tại Chương~\ref{chap:phantich}, hệ thống Farmery được hiện thực hóa thông qua 8 mô-đun nghiệp vụ trọng tâm. Mỗi mô-đun được phân rã thành các dịch vụ độc lập trên máy chủ nghiệp vụ NestJS, giao tiếp qua bus sự kiện nội bộ và ràng buộc dữ liệu chặt chẽ theo hợp đồng giao diện OpenAPI.

\subsubsection*{Module xác thực \& quản lý tài khoản}
Mô-đun tài khoản (\texttt{identity}) đảm nhiệm vai trò định danh, cấp phát quyền hạn và bảo vệ an toàn giao dịch cho toàn bộ các tác nhân trong hệ thống:
\begin{itemize}
    \item \textbf{Xác thực không mật khẩu (Passwordless Authentication):} Người dùng đăng nhập thông qua mã xác thực một lần (OTP) 6 chữ số gửi trực tiếp qua email. Nhằm ngăn chặn lộ lọt thông tin, máy chủ chỉ lưu trữ giá trị băm HMAC của mã OTP trong bộ nhớ tạm thời với thời gian sống (TTL) 10 phút, giới hạn tối đa 5 lần thử sai và áp dụng cơ chế điều tiết tần suất gửi mã (\emph{rate limiting}) để chống tấn công vét cạn.
    \item \textbf{Phân quyền dựa trên vai trò (RBAC):} Hệ thống phân định nghiêm ngặt 6 nhóm vai trò bao gồm khách lẻ (\texttt{buyer\_retail}), doanh nghiệp mua sỉ (\texttt{buyer\_wholesale}), nhà cung cấp (\texttt{seller}), cán bộ kiểm duyệt (\texttt{staff\_moderation}), nhân viên kho bãi/vận hành (\texttt{staff\_warehouse}, \texttt{staff\_operations}) và chủ nền tảng (\texttt{platform\_owner}). Quyền hạn được đóng gói trong mã thông báo JWT ký điện tử, được chặn và kiểm tra tự động tại tầng cổng bảo vệ (\texttt{AuthGuard}).
\end{itemize}

\subsubsection*{Module sản phẩm \& truy xuất nguồn gốc (QR)}
Mô-đun danh mục và nguồn cung (\texttt{catalog}) là hạt nhân đảm bảo tính minh bạch xuất xứ nông sản từ nông trại đến người tiêu dùng:
\begin{itemize}
    \item \textbf{Quản lý lô hàng và nhật ký canh tác:} Nhà cung cấp khai báo lô hàng gắn liền với vùng trồng và thời vụ cụ thể. Hệ thống hỗ trợ sản lượng lô dưới dạng số thập phân (ví dụ 0,5~kg đối với các loại nông sản giá trị cao như nấm linh chi, dâu tây giống VIP), chuẩn hóa toàn bộ mốc thời gian theo giờ Việt Nam (\texttt{Asia/Ho\_Chi\_Minh}). Đồng thời, API thực hiện kiểm tra chặt chẽ quyền sở hữu, đảm bảo nhà cung cấp chỉ được phép tạo lô hàng cho các sản phẩm đã được phê duyệt của chính đơn vị mình.
    \item \textbf{Nhật ký canh tác bất biến và cơ chế đính chính:} Nhằm chống gian lận dữ liệu nông nghiệp, các dòng nhật ký canh tác được tổ chức theo mô hình chỉ thêm (\emph{append-only}). Khi cần sửa đổi thông tin sai sót, nhà cung cấp thực hiện đính chính bằng cách tạo một bản ghi mới liên kết tường minh qua trường \texttt{supersedesId}; hệ thống tự động ẩn dòng cũ, hiển thị dòng đính chính mới nhất và ngăn chặn việc đính chính nhiều lần trên cùng một dòng nhật ký gốc.
    \item \textbf{Hỗ trợ ghi nhật ký ngoại tuyến và mẫu theo loại cây trồng:} Giao diện khu nhà cung cấp tích hợp cơ chế tự động sao lưu bản thảo nhật ký vào bộ nhớ cục bộ trình duyệt (\texttt{localStorage}) khi mất kết nối mạng, giúp nông dân không bị mất dữ liệu khi đang thao tác ngoài đồng ruộng. Ngoài ra, biểu mẫu gợi ý sẵn các mẫu thao tác canh tác chuẩn VietGAP chuyên biệt cho từng nhóm cây trồng: Nấm, Rau ăn lá, Cây ăn trái, Dâu tây và Cà phê/Trà.
    \item \textbf{Tem QR động và trang truy xuất nguồn gốc:} Sau khi lô hàng được kiểm duyệt viên phê duyệt, hệ thống kích hoạt mã QR định danh duy nhất dẫn tới trang truy xuất công khai (\texttt{/truy-xuat/[ma]}). Giao diện cung cấp bố cục in ấn chuyên dụng (\texttt{\#tem-qr-print}) tối ưu cho máy in nhiệt dán tem nhãn bao bì và tính năng xuất tệp ảnh QR độ phân giải cao phục vụ truyền thông, chia sẻ.
    \item \textbf{Hồ sơ nhà cung cấp công khai:} Xây dựng trang hồ sơ định danh điện tử của từng nhà cung cấp (\texttt{/ho-so-nha-cung-cap/[id]}), liên kết trực tiếp từ thông tin vùng trồng ở trang chi tiết sản phẩm, công khai đầy đủ chứng nhận chất lượng (VietGAP, GlobalGAP, OCOP) đã được cơ quan chức năng hậu kiểm.
\end{itemize}

\subsubsection*{Module bán lẻ (B2C)}
Mô-đun bán lẻ (\texttt{retail}) phục vụ luồng mua sắm của khách hàng cá nhân với cam kết nông sản tươi ngon giao nhanh:
\begin{itemize}
    \item \textbf{Khung mua hàng thông minh:} Tự động phân loại quy cách đóng gói, cho phép nhập số lượng lẻ đối với nông sản bán theo khối lượng (kg) và bắt buộc số nguyên đối với hàng đóng khay/hộp.
    \item \textbf{Khóa dòng tồn và cơ chế chống bán vượt (Anti-overselling):} Để triệt tiêu tình trạng nhiều khách hàng cùng tranh chấp phần tồn kho cuối cùng của một lô hàng tươi sống, hệ thống áp dụng cơ chế khóa dòng bi quan (\emph{pessimistic locking}) trực tiếp bằng câu lệnh SQL thô \texttt{SELECT ... FOR UPDATE} trên bảng dòng tồn. Phần hàng được tạm giữ chỗ trong vòng 15 phút để khách hoàn tất thanh toán; quá thời hạn, hệ thống tự động giải phóng tồn kho về trạng thái khả dụng.
    \item \textbf{Thanh toán và COD có điều kiện:} Hỗ trợ thanh toán điện tử qua cổng giả lập và phương thức trả tiền khi nhận hàng (COD). Nhằm bảo vệ quyền lợi của nông dân và nền tảng trước nguy cơ hủy đơn ảo, hệ thống tự động kiểm tra lịch sử từ chối nhận hàng để từ chối hình thức COD đối với các tài khoản từng vi phạm.
\end{itemize}

\subsubsection*{Module đánh giá \& uy tín}
Mô-đun đánh giá (\texttt{review}) thiết lập cơ chế phản hồi hai chiều minh bạch giữa người mua và nhà sản xuất:
\begin{itemize}
    \item \textbf{Thu thập đánh giá gắn với đơn hàng thực tế:} Chỉ những tài khoản đã nhận hàng thành công mới được quyền gửi đánh giá sao và nhận xét chi tiết kèm hình ảnh cho từng sản phẩm của lô hàng tương ứng.
    \item \textbf{Kiểm duyệt và tính điểm uy tín tổng hợp:} Đánh giá phải trải qua hàng đợi kiểm duyệt tự động kết hợp thủ công để loại trừ ngôn từ vi phạm tiêu chuẩn cộng đồng. Điểm uy tín của nhà cung cấp được tính toán tự động dựa trên trọng số của điểm đánh giá người mua, tỷ lệ hoàn tất đơn hàng và tỷ lệ phát sinh khiếu nại chất lượng.
\end{itemize}

\subsubsection*{Module bán sỉ (B2B)}
Mô-đun bán sỉ (\texttt{wholesale}) tối ưu hóa quy trình phân phối sản lượng lớn giữa nông hộ/hợp tác xã và các doanh nghiệp thu mua:
\begin{itemize}
    \item \textbf{Bảng giá sỉ theo bậc (Tiered Pricing):} Nhà cung cấp thiết lập các bậc giá tương ứng với ngưỡng sản lượng tối thiểu (MOQ) và thời gian chuẩn bị hàng. Khung mua hàng phía người mua tự động nhận diện sản lượng đặt để áp dụng mức đơn giá chiết khấu tối ưu nhất.
    \item \textbf{Thương lượng một vòng phản hồi giá (Single-round RFQ):} Nhằm đơn giản hóa thủ tục và giảm thiểu độ trễ giao dịch nông sản tươi, hệ thống quy định cơ chế thương lượng rút gọn duy nhất một vòng. Doanh nghiệp mua gửi yêu cầu báo giá (RFQ) $\rightarrow$ Nhà cung cấp gửi lại báo giá phản hồi duy nhất một lượt (chốt đơn giá, sản lượng đáp ứng và ngày giao cam kết) $\rightarrow$ Doanh nghiệp mua đưa ra quyết định cuối cùng: chấp nhận hoặc từ chối để đóng yêu cầu.
    \item \textbf{Hợp đồng điện tử mẫu kèm mã băm SHA-256:} Khi bên mua chấp thuận mức giá phản hồi, hệ thống lập tức tự sinh một hợp đồng mua bán nông sản sỉ mẫu chuẩn (\texttt{Contract}) với mã số duy nhất (ví dụ \texttt{HD-2026-8002}). Văn bản hợp đồng được băm bằng thuật toán SHA-256 (\texttt{documentHash}) lưu trữ đồng hành để bảo đảm tính toàn vẹn và chống chối bỏ nội dung thỏa thuận.
    \item \textbf{Xác nhận giao kết qua OTP và thanh toán giả lập:} Đại diện doanh nghiệp mua ký xác nhận giao kết hợp đồng thông qua mã OTP gửi về email định danh. Sau khi hợp đồng được giao kết (\texttt{concluded}), hệ thống chuyển sang bước thanh toán sỉ thông qua dịch vụ thanh toán \texttt{PaymentService} (sử dụng cổng thanh toán giả lập \texttt{DemoPaymentGateway}) và thiết lập lịch trình giao hàng một đợt (\texttt{single\_batch}) tập trung tại kho trung chuyển Farmery.
\end{itemize}

\subsubsection*{Module quản lý tồn kho \& vận chuyển/khiếu nại}
Mô-đun kho vận (\texttt{inventory} \& \texttt{fulfillment}) đảm bảo tính liên tục của chuỗi cung ứng nông sản ba chặng:
\begin{itemize}
    \item \textbf{Quản lý tồn kho theo lô (FEFO):} Toàn bộ hàng hóa nhập kho được quản lý định danh theo từng lô thu hoạch. Thuật toán xuất kho ưu tiên lô có hạn sử dụng gần nhất xuất trước (\emph{First Expired, First Out}) nhằm giảm thiểu tỷ lệ hao hụt nông sản.
    \item \textbf{Quản lý đơn thu mua phía nhà cung cấp (FR13.1):} Giao diện khu nhà cung cấp cho phép theo dõi các đơn thu mua từ nền tảng, hiển thị tách biệt sản lượng theo từng quy cách đóng gói (không cộng gộp lẫn lộn giữa kg và hộp) và giám sát tỷ lệ nghiệm thu nhập kho thực tế của nhân viên kiểm hàng.
    \item \textbf{Điều phối vận chuyển ba chặng:} Quản lý lộ trình gom hàng từ trang trại về kho sơ chế vùng (Chặng~1), lập chuyến xe trung chuyển container lạnh giữa các kho trung tâm (Chặng~2) và phân bổ đơn hàng cho nhân viên giao hàng chặng cuối (Chặng~3).
\end{itemize}

\subsubsection*{Module quản trị}
Mô-đun quản trị (\texttt{platform}) cung cấp công cụ vận hành toàn diện cho nhân sự nội bộ thông qua ứng dụng web chuyên biệt:
\begin{itemize}
    \item \textbf{Hàng đợi kiểm duyệt đa cấp:} Cán bộ kiểm duyệt tiếp nhận và phê duyệt hồ sơ năng lực nhà cung cấp, bản quét chứng nhận an toàn thực phẩm, mẫu đăng ký sản phẩm mới và hồ sơ lô hàng xuất chuồng/thu hoạch trước khi cấp quyền lưu hành mã QR.
    \item \textbf{Giám sát nghiệp vụ và xử lý vi phạm:} Cung cấp bảng điều khiển trực quan theo dõi sản lượng bán lẻ, biến động giá bán sỉ, cảnh báo thu hồi lô hàng khi phát hiện mẫu kiểm nghiệm không đạt tiêu chuẩn và công cụ tạm khóa tài khoản khi có dấu hiệu gian lận thương mại.
\end{itemize}

\subsubsection*{Module AI}
\label{sec:hienthucAI}
Mô-đun trí tuệ nhân tạo được xây dựng dưới dạng dịch vụ vi mô độc lập bằng Python (FastAPI), tích hợp sâu vào hệ thống qua các cổng API:
\begin{itemize}
    \item \textbf{Tìm kiếm thông minh và gợi ý tự động (Search Suggestion):} Hệ thống phân tích cú pháp câu truy vấn người dùng trong thời gian thực, tự động bóc tách từ khóa sản phẩm khỏi các điều kiện lọc (chứng nhận VietGAP/hữu cơ, khoảng giá bán lẻ, vùng trồng Đà Lạt/Đơn Dương). Bộ gợi ý xử lý trơn tru các trường hợp gõ tiếng Việt không dấu, viết dính từ và sai lỗi chính tả phổ biến.
    \item \textbf{Tìm kiếm ngữ nghĩa lai ghép (Hybrid Search):} Ứng dụng mô hình nhúng đa ngôn ngữ tiên tiến BGE-M3 để chuyển đổi mô tả sản phẩm và nhu cầu tìm kiếm sang không gian vector 1024 chiều lưu trữ trên \texttt{pgvector}. Kết quả vector được kết hợp với thuật toán tìm kiếm từ khóa truyền thống thông qua giải thuật dung hợp thứ hạng tương hỗ RRF (\emph{Reciprocal Rank Fusion}), vừa hiểu được ý định trừu tượng của người mua (ví dụ \lq\lq rau ăn lẩu\rq\rq, \lq\lq hoa quả làm quà biếu\rq\rq), vừa bảo đảm tính chính xác của tên sản phẩm riêng biệt.
    \item \textbf{Dự báo biến động giá nông sản:} Thu thập chuỗi dữ liệu giá công khai từ các bản tin thị trường định kỳ của Sở Công Thương và Cục Thống kê, trích xuất dữ liệu chuẩn hóa dạng bảng (CSV) theo tuần và vùng trồng. Các mô hình chuỗi thời gian SARIMA và Prophet được huấn luyện để đưa ra biên độ giá tham khảo cho nhà cung cấp khi thiết lập bảng giá bậc, giúp bình ổn giá và dự báo xu hướng cung cầu theo mùa vụ.
\end{itemize}

\subsection{Giải thuật và đoạn code quan trọng}
\label{sec:giaithuat}
% TODO: trích các đoạn code/giải thuật tiêu biểu (vd sinh mã QR, tính giá bậc theo MOQ, middleware RBAC) kèm giải thích ngắn gọn.


\section{Kiểm thử}
\label{sec:ketquachucnang}
% TODO: kết quả thực thi các test case thiết kế tại mục~\ref{sec:kichbankiemthu} (đạt/không đạt, lỗi phát hiện và cách khắc phục).

\section{Kết chương}
% TODO
